\chapter{Gauge Theory Taxonomy For Deforming Bodies}
\label{app:gauge-taxonomy}

The deforming-body chapter uses gauge language because a body can change shape
while a symmetry group moves it through space.  This appendix separates the
main connection concepts that appear in rigid-body reconstruction, swimming,
elastic frames, and Cosserat media.  The common pattern is a split between
``where the body is'' and ``what shape the body has.''

\section{Principal-Bundle Kinematics}
\label{sec:gauge-principal-bundle-kinematics}

Let \(Q\) be the configuration space of a body and let a Lie group \(G\) act
freely on \(Q\).  For a rigid body in space, \(G=\SE(3)\) or \(G=\SO(3)\).
For a planar swimmer, \(G=\SE(2)\).  The quotient
\[
        S=Q/G
\]
is the shape space.  Locally one may write a configuration as
\[
        q=(g,s),\qquad g\in G,\quad s\in S.
\]

\paragraph{Local notation.}
In this appendix \(Q\) is a configuration space, \(G\) a symmetry group,
\(S\) a shape space, \(g\) a group element, \(s\) a shape coordinate, and
\(\xi=g^{-1}\dot g\) a body velocity.  The connection one-form \(A\) in this
appendix is gauge data, not the Laplace-Runge-Lenz vector and not an
electromagnetic vector potential.

The body velocity is
\[
        \xi=g^{-1}\dot g\in\mathfrak{g}
\]
for a left-trivialized convention.  A connection is a rule that assigns a
group velocity to a shape velocity:
\[
        \xi + A(s)\dot s =0
\]
for zero-momentum or force-free motion.  The sign is conventional; what matters
is that \(A\) is a \(\mathfrak{g}\)-valued one-form on shape space.

\begin{figure}[htbp]
\centering
\begin{tikzpicture}[>=Stealth, scale=1.0]
\draw[thick] (-3.0,-0.8) rectangle (3.0,1.0);
\draw[thick] (-2.2,-2.0) rectangle (2.2,-1.2);
\node at (0,1.25) {configuration space \(Q\)};
\node at (0,-2.3) {shape space \(S=Q/G\)};
\draw[->, thick] (0,-0.8) -- (0,-1.2) node[midway, right] {\(\pi\)};
\draw[thick, blue] (-2.3,0.05) .. controls (-1.2,0.8) and (1.1,0.8) .. (2.2,0.05);
\draw[thick, red] (-1.5,-1.62) .. controls (-0.5,-1.35) and (0.6,-1.9) .. (1.5,-1.62);
\draw[->, thick] (-1.1,-1.45) -- (-0.4,-1.43);
\draw[->, thick, blue] (-0.4,0.53) -- (0.35,0.55);
\node[blue] at (2.55,0.35) {horizontal lift};
\node[red] at (1.85,-1.42) {shape path};
\draw[->] (-2.45,0.1) -- (-2.45,0.7) node[midway,left] {\(G\)-fiber};
\end{tikzpicture}
\caption{A connection chooses horizontal lifts of shape-space motion.  The
failure of a lifted loop to close in \(Q\) is the geometric phase or holonomy.}
\label{fig:gauge-principal-bundle}
\end{figure}

\section{Mechanical Connection}
\label{sec:mechanical-connection-taxonomy}

Suppose the kinetic energy defines a metric on \(Q\) and the \(G\)-action is a
symmetry.  The momentum map \(J:TQ\to\mathfrak{g}^*\) measures motion along the
group directions.  The mechanical connection is the projection that declares a
velocity horizontal when its group momentum vanishes:
\[
        J(q,\dot q)=0.
\]
In local variables \((g,s)\), the kinetic energy has the schematic form
\[
        T
        =
        \frac{1}{2}
        \begin{pmatrix}\xi\\ \dot s\end{pmatrix}^T
        \begin{pmatrix}
        \mathcal{I}(s) & \mathcal{K}(s)\\
        \mathcal{K}(s)^T & \mathcal{M}(s)
        \end{pmatrix}
        \begin{pmatrix}\xi\\ \dot s\end{pmatrix}.
\]
The group momentum is
\[
        \Pi=\frac{\pd T}{\pd \xi}
        =
        \mathcal{I}(s)\xi+\mathcal{K}(s)\dot s .
\]
At zero momentum,
\[
        \xi=-\mathcal{I}(s)^{-1}\mathcal{K}(s)\dot s.
\]
Thus
\[
        A(s)=\mathcal{I}(s)^{-1}\mathcal{K}(s).
\]
This is the inertial connection used for falling cats, deforming rigid
assemblies, and the finite-dimensional examples in the main text.

\section{Local Connection For Swimming}
\label{sec:local-connection-swimming}

At low Reynolds number, inertia is negligible.  The governing equations are
force and torque balance with a linear relation between body velocity and
hydrodynamic force.  In local coordinates one again obtains
\[
        \xi=-A_{\rm Stokes}(s)\dot s,
\]
but the origin of \(A_{\rm Stokes}\) is not kinetic energy.  It is the
resistance problem for the surrounding fluid.  The same gauge geometry appears
because group motion is reconstructed from shape motion, but the physical
inner product and the equations that define horizontality are different.

\section{Gauge Transformation}
\label{sec:gauge-transformation-taxonomy}

A local trivialization \(q=(g,s)\) is a choice of body frame over shape space.
Changing the frame by a smooth map \(h:S\to G\) gives
\[
        g'=g h(s).
\]
For the left-trivialized velocity \(\xi=g^{-1}\dot g\), the connection
one-form transforms as
\[
        A'
        =
        \Ad_{h^{-1}}A-h^{-1}\dd h .
\]
The minus sign follows from the reconstruction convention
\(\xi+A(s)\dot s=0\). If one instead defines the standard connection
\(B=-A\), then \(B'=\Ad_{h^{-1}}B+h^{-1}\dd h\). This formula explains the
word ``gauge.''  The local coefficients of connection depend on the chosen
frame, but geometric predictions do not. Holonomy around a closed shape loop
changes by conjugation, and conjugacy invariants are frame-independent.

\section{Curvature And Small Loops}
\label{sec:gauge-curvature-taxonomy}

For the reconstruction convention used in the finite-dimensional
deforming-body chapter, the curvature is the \(\mathfrak{g}\)-valued two-form
\[
        F_A=\dd A-\frac{1}{2}[A\wedge A].
\]
The bracketed wedge is defined so that
\([A\wedge A]_{\alpha\beta}=2[A_\alpha,A_\beta]\). Equivalently, the component
formula below is the convention being used. This is the same convention used in
Chapter~\ref{chap:deforming-body} for
\(\mathcal F=\dd\mathcal A-\frac12[\mathcal A\wedge\mathcal A]\); the factor
\(\frac12\) is the wedge-symmetrization factor, not a change of sign or a
different curvature.
For a small oriented loop \(\Sigma\) in shape space, the leading geometric
phase is
\[
        g_{\rm hol}
        =
        \exp\left(-\int_\Sigma F_A+O(\operatorname{area}^{3/2})\right)
\]
when noncommutativity is weak over the loop.  In two shape variables
\((\alpha,\beta)\), if
\[
        A=A_\alpha(\alpha,\beta)\,\dd\alpha+
          A_\beta(\alpha,\beta)\,\dd\beta,
\]
then
\[
        F_{\alpha\beta}
        =
        \pd_\alpha A_\beta-\pd_\beta A_\alpha-
        [A_\alpha,A_\beta].
\]
The bracket term is the non-Abelian correction in this reconstruction sign
convention.  It vanishes for a translation group \(\R^m\), but not for
rotations or rigid motions.

\section{Connection Types In These Notes}
\label{sec:gauge-connection-types-table}

\begin{center}
\begin{tabular}{p{0.26\textwidth} p{0.30\textwidth} p{0.34\textwidth}}
\hline
connection & horizontality condition & typical use\\
\hline
mechanical connection & zero group momentum \(J=0\) & falling cat, internal
rotors, deforming inertial bodies\\
Stokes or local swimming connection & zero net force and torque in viscous
flow & low-Reynolds swimming and crawling\\
material frame connection & comparison of frames attached to a continuum &
elastic directors, rods, shells\\
spin connection & covariant differentiation of orthonormal frames & curved
surfaces and continuum frame fields\\
Cosserat connection & director rotation along a material coordinate & rods and
micropolar continua\\
defect connection & incompatible local frames or coframes & disclinations,
dislocations, torsion, curvature\\
\hline
\end{tabular}
\end{center}

These objects share gauge transformation laws, but they are not the same
physical connection.  The defining horizontal condition must always be stated.

\section{Relation To Elasticity And Fluids}
\label{sec:gauge-relation-elastic-fluid}

Elasticity uses a deformation gradient \(F\) and may also use local directors.
Gauge language enters when one compares material frames across the body or
separates a rigid group motion from internal shape.  A Cosserat rod, for
example, has a centerline \(r(s)\) and an orthonormal frame
\((d_1(s),d_2(s),d_3(s))\).  The strain vector \(\kappa(s)\) is defined by
\[
        \pd_s d_a=\kappa\times d_a .
\]
This is a one-dimensional connection along the rod.

Fluids use the diffeomorphism group more directly.  In ideal fluid mechanics,
particle relabeling symmetry leads to Kelvin circulation and the Euler
equation.  In viscous fluid mechanics, dissipation breaks the purely
Hamiltonian picture, but geometric structures remain useful for vorticity,
transport, and reduction.  The gauge-theory chapter focuses on deforming
bodies because the finite-dimensional connection picture can be displayed
there with the least analytic overhead.

\section{Fluid Relabeling Versus Shape-Space Gauge Theory}
\label{sec:fluid-relabeling-versus-shape-gauge}

Fluid mechanics has a gauge-like symmetry that is easy to confuse with the
Shapere-Wilczek mechanism.  In a material description, a fluid motion is a
time-dependent diffeomorphism
\[
        \eta_t:\mathcal{B}\to\Omega .
\]
Relabeling particles by a volume-preserving diffeomorphism
\(\psi:\mathcal{B}\to\mathcal{B}\) changes \(\eta_t\) to
\(\eta_t\circ\psi\) without changing the Eulerian velocity field
\[
        v(x,t)=\dot\eta_t(\eta_t^{-1}(x)).
\]
The quotient by particle relabeling is therefore the passage from material
labels to Eulerian fields.  Kelvin's circulation theorem can be understood as
a consequence of this relabeling symmetry in the ideal-fluid variational
principle.

For a deforming swimmer or falling cat, the quotient has a different meaning.
The base variables are shape variables \(s\), the fiber is a physical rigid
motion \(g\in G\), and the connection reconstructs \(g(t)\) from a path in
shape space.  The holonomy is an observable displacement or rotation.  Fluid
relabeling removes unobservable labels; shape-space gauge theory reconstructs
observable group motion from internal deformation.  Both use quotient
geometry, but they answer different physical questions.

\section{Common Pitfalls}
\label{sec:gauge-common-pitfalls}

\begin{enumerate}
\item Do not identify gauge dependence with physical ambiguity.  Coordinates
and frames are choices; a chosen frame gives a representative holonomy, while
frame-independent statements are its conjugacy-invariant content and the
reconstructed motion relative to specified boundary frames.
\item Do not use the mechanical connection for a low-Reynolds swimmer unless
the kinetic-energy metric is actually the relevant physics.  The Stokes
connection comes from resistance, not inertia.
\item Do not infer zero phase from zero instantaneous group velocity in one
frame.  A closed shape loop can produce net motion through curvature.
\item Do not ignore noncommutativity for rotations.  For \(G=\SO(3)\) or
\(\SE(3)\), the bracket term in \(F_A\) and the ordering of finite motions
matter.
\item Do not call every internal coordinate a gauge variable.  The gauge group
is a redundancy or symmetry acting along fibers; shape variables label the
base.
\end{enumerate}

The conceptual payoff is that the rigid body, deforming body, elastic director
field, and swimmer are not isolated examples.  They are different realizations
of a common reconstruction problem: shape motion plus a connection determines
motion in a symmetry group.
